\documentclass[../../main.tex]{subfiles}

\begin{document}

\subsection{Il tracciatore} \label{subsec:il-tracciatore-interno}
  \begin{figure}[ht]
    \centering
    \includegraphics[width=.5\textwidth]{DUNE/sand/stt-section}
    \caption{
      Schema di un modulo STT comprendente i tre elementi principali (
      da sinistra a destra): un bersaglio di polipropilene \ce{CH2} regolabile,
      un radiatore con 119 fogli di polipropilene per l'identificazione di
      $\Pepm$ e quattro strati di \textit{straw} in configurazione XXYY (fascio
      lungo l'asse z e campo B lungo l'asse x).
      Alcuni dei fogli del radiatore sono sostituiti da sottili bersagli di
      carbonio. \cite{abud2021deepundergroundneutrinoexperiment}
    }
    \label{fig:stt-section}
  \end{figure}
  Il tracciatore, che occupa la maggior parte del volume all'interno dell'ECAL,
  ha due obiettivi primari: offrire un bersaglio per le interazioni di
  (anti)neutrini e fornire un sistema di tracciamento in grado di ricostruire
  accuratamente il momento di tutte le particelle cariche.
  Per raggiungere questi obiettivi sono stati identificati alcuni
  requisiti fisici fondamentali del tracciatore:
  \begin{itemize}
    \item avere una bassa densità e un alto campionamento delle tracce per
    ridurre al minimo lo scattering multiplo e garantire così un'elevata
    risoluzione spaziale e della quantità di moto;

    \item avere uno spessore totale confrontabile con la lunghezza
    di radiazione per minimizzare le interazioni secondarie;

    \item essere in grado di ospitare diversi materiali bersaglio per
    confrontare le interazioni dei neutrini su differenti nuclei all'interno
    dello stesso rivelatore;

    \item essere in grado di distinguere $\Pepm, \Ppipm, \PKpm, \Pproton$
    e $\Pmupm$;

    \item avere una massa fiduciale del bersaglio di $\sim$5 tonnellate per
    garantire un tasso di interazione sufficiente.
  \end{itemize}
  Lo STT di SAND segue un design modulare; ogni modulo, il cui schema è
  illustrato in figura ~\ref{fig:stt-section}, è equipaggiato con:
  \begin{itemize}
    \item uno strato di bersaglio passivo in carbonio o polipropilene con uno
    spessore rispettivamente di 4~mm o 5~mm;

    \item un radiatore composto da 105 fogli di polipropilene spessi
    \SI{18}{\micro\meter}, alternati a intercapedini d'aria di
    \SI{117}{\micro\meter} per separare $\Pepm$ da $\Ppipm$
    tramite radiazione di transizione;

    \item quattro strati di \textit{straw}, dove ogni cannuccia ha una lunghezza
    di 3,2~m, un diametro di 5~mm, pareti in Mylar da \SI{12}{\micro\meter}
    rivestite di alluminio e un filo di tungsteno da \SI{20}{\micro\meter}
    rivestito di oro, opera con una miscela gassosa di \ce{Xe} e \ce{CO2}
    (70/30) a 1,9~atm; due strati di straw sono disposti lungo l'asse X e gli
    altri due lungo l'asse Y in una configurazione XXYY\@.
  \end{itemize}
  Gli strati del bersaglio e del radiatore costituiscono i bersagli principali
  del rivelatore, al posto delle lastre bersaglio è possibile installare molti
  materiali diversi come \ce{C}, \ce{Ca}, \ce{Fe} e \ce{Pb}.\\
  La densità media del rivelatore può essere modulata tra un massimo di
  \SI{0,18}{g/cm^3}, corrispondente allo spessore del radiatore e del
  bersaglio, e un minimo di \SI{0,005}{g/cm^3} se sono presenti solo gli strati
  di \textit{straw}.
  Questa densità è mantenuta sufficientemente bassa per ottenere una lunghezza
  totale del rivelatore confrontabile con una lunghezza di radiazione,
  consentendo una misura accurata dei quadrimomenti delle particelle finali.\\
  I bersagli passivi costituiscono fino al 97\% della massa totale del
  rivelatore, mentre il restante 3\% è dovuto alla massa delle \textit{straw}.\\
  Lo STT fornisce: una risoluzione spaziale del singolo punto inferiore ai
  \SI{200}{\micro\meter}, una risoluzione dell'impulso di $\sim$3\% e una
  risoluzione angolare di 2~mrad.

\end{document}